\section{Universal norms and one-prime ramification at seven}
\label{seven:universal-norms}

The finite strict module of Theorem~\ref{seven:main} does not coincide
with the Iwasawa module in which ramification is allowed at one of the
two primes above seven. We determine the latter module and its unit
quotient. The distinction comes from capitulation: the faithful class
groups remain of order seven under norms, but their extension maps
vanish, and the resulting defect in norms of global units persists
throughout the tower.

Fix a prime $\mathfrak p\mid7$ of $K=\mathbb Q(\omega)$ and a faithful
character $\chi$ of $\Delta=\operatorname{Gal}(E/K)$. Choose
$\gamma\in\Gamma=\operatorname{Gal}(E_\infty/E)$ with
$\langle\varepsilon_7\rangle(\gamma)=8$, and put
\[
 \Lambda=\mathbb Z_7[[T]],\qquad T=\gamma-1,\qquad
 \omega_n(T)=(1+T)^{7^n}-1.
\]
Let $L_n$ be the completion of $E_n$ at its unique prime over
$\mathfrak p$, and define
\[
 U_n=U^1(L_n),\qquad
 \mathcal E_n=\mathbb Z_7\otimes_{\mathbb Z}\mathcal O_{E_n}^{\times},
 \qquad
 U_\infty=\varprojlim_N U_n,\qquad
 \mathcal E_\infty=\varprojlim_N\mathcal E_n.
\]
The localization map $\ell_n:\mathcal E_n\to U_n$ uses the principal-unit
projection. Its image $\overline{\mathcal E}_n$ is the closure of the
global units which are principal at this prime. Write
$\overline{\mathcal E}_\infty=\varprojlim_N\overline{\mathcal E}_n$ and
\[
 V_{\mathfrak p,\chi}
   =(U_\infty/\overline{\mathcal E}_\infty)^\chi.
\]
Finally, let $X_{\mathfrak p,\infty}$ be the Galois group of the maximal
abelian pro-seven extension of $E_\infty$ unramified outside the primes
over $\mathfrak p$. These are the one-prime local-unit and Galois modules
of \cite[printed p.~36]{Rubin}.

\begin{theorem}\label{seven:one-prime-linear}
For each choice of $\mathfrak p$ and faithful $\chi$, there is
$a_{\mathfrak p,\chi}\in7\mathbb Z_7^\times$ such that
\[
 U_\infty^\chi\simeq\Lambda,\qquad
 \mathcal E_\infty^\chi\simeq\Lambda,\qquad
 V_{\mathfrak p,\chi}\simeq
 X_{\mathfrak p,\infty}^\chi
 \simeq\Lambda/(T-a_{\mathfrak p,\chi}).
\]
The class-field exact sequence is
\begin{equation*}\tag{U1}\label{seven:one-prime-exact}
 0\longrightarrow V_{\mathfrak p,\chi}
 \longrightarrow X_{\mathfrak p,\infty}^\chi
 \longrightarrow M^\chi
 \longrightarrow0,
\end{equation*}
where $M^\chi=\Lambda/(7,T)$, and its first map identifies
$V_{\mathfrak p,\chi}$ with $7X_{\mathfrak p,\infty}^\chi$.
In particular the one-prime characteristic ideal has
$\mu=0$, $\lambda=1$, and is not the unit ideal.

Put
\[
 e_{\mathfrak p,\chi}^{*}
 =-\frac{\log(1+a_{\mathfrak p,\chi})}{\log8}
 \in\mathbb Z_7^\times.
\]
For every generator $f_{\mathfrak p,\chi}$ of
$\operatorname{char}_{\Lambda}X_{\mathfrak p,\infty}^\chi$ and every
$e\in\mathbb Z_7$,
\begin{equation*}\tag{U2}\label{seven:one-prime-weight-valuation}
 v_7\!\left(f_{\mathfrak p,\chi}(8^{-e}-1)\right)
 =1+v_7(e-e_{\mathfrak p,\chi}^{*}),
\end{equation*}
with both sides interpreted as infinity at
$e=e_{\mathfrak p,\chi}^{*}$. These valuations are unbounded on the
positive integer weights, whereas their strict Selmer lengths in
Theorem~\ref{seven:main} are all one.
\end{theorem}

\subsection{Capitulation and norms of global units}

All module notation for completed units in this proof is additive.
Theorem~\ref{seven:main} gives
$A(E_n)^\chi\simeq C_7$ and norm isomorphisms between these groups.
For $n>m$, extension and norm of ideals satisfy
\[
 N_{n/m}j_{n/m}=7^{n-m}
 \quad\hbox{on } A(E_m)^\chi.
\]
Since the right side vanishes and the norm is injective, we have
\begin{equation*}\tag{U3}\label{seven:faithful-capitulation}
 j_{n/m}=0.
\end{equation*}

\begin{lemma}\label{seven:unit-norm-index}
For every $n>m$,
\[
 H^1(\operatorname{Gal}(E_n/E_m),\mathcal E_n)^\chi\simeq C_7,
 \qquad
 [\mathcal E_m^\chi:N_{n/m}\mathcal E_n^\chi]=7.
\]
In particular
$N_{n/0}\mathcal E_n^\chi=7\mathcal E_0^\chi$ for $n\ge1$.
\end{lemma}

\begin{proof}
Set $L=E_n$, $k=E_m$, and $G=\operatorname{Gal}(L/k)$. Hilbert 90 and
the principal-ideal sequence identify
$H^1(G,\mathcal O_L^\times)$ with $P_L^G/P_k$.
The natural map
\[
 P_L^G/P_k\longrightarrow I_L^G/jI_k
 \simeq
 \bigoplus_{\mathfrak q\ {\rm ramified\ in}\ L/k}
           \mathbb Z/e_{\mathfrak q}\mathbb Z
\]
has kernel equal to the capitulation kernel on ideal classes.
The only ramified primes are the two primes above seven, and each is
fixed by $\Delta$. The displayed direct sum is therefore
$\Delta$-trivial. Applying a faithful idempotent and then completing
at seven proves
\[
 H^1(G,\mathcal E_n)^\chi
 \simeq\ker(j_{n/m})^\chi\simeq C_7.
\]
Completion commutes with this cohomology because $\mathbb Z_7$ is flat
over $\mathbb Z$ and the group is finite.

The complex places of $E_n$ form a regular set for
$\Delta\times\operatorname{Gal}(E_n/E)$. Indeed, each complex place has
exactly one embedding extending a fixed embedding of the imaginary
quadratic field $K$. Dirichlet's unit theorem consequently gives the
regular representation with one trivial representation removed.
On a faithful character the removed term disappears, so
\[
 \mathbb Q_7\otimes_{\mathbb Z_7}\mathcal E_n^\chi
 \simeq
 \mathbb Q_7[\operatorname{Gal}(E_n/E)].
\]
Its restriction to $G$ is a sum of regular representations. Hence the
Herbrand quotient of $\mathcal E_n^\chi$ is one: an integral regular
lattice has trivial Tate cohomology, and replacing a lattice by a
commensurable one does not change its Herbrand quotient, since finite
modules have quotient one. Thus
\[
 \#\widehat H^0(G,\mathcal E_n^\chi)
 =\#H^1(G,\mathcal E_n^\chi)=7.
\]
The invariant units are $\mathcal E_m^\chi$, proving the asserted
norm index. At level zero $\mathcal E_0^\chi$ has rank one, so its
submodule of index seven is $7\mathcal E_0^\chi$.
\end{proof}

We record explicitly the base localization. Put
\[
 \varepsilon_0=-270+739\alpha-108\alpha^2.
\]
Reduction using $\alpha^3=269$ gives
\[
 \varepsilon_0^{342}
 \equiv1+7\alpha+28\alpha^2\pmod{49}.
\]
Since $342\equiv6\pmod7$, the principal logarithm therefore satisfies
\begin{equation*}\tag{U4}\label{seven:base-local-generator}
 \frac{\log\varepsilon_0}{7}
 \equiv6\alpha+3\alpha^2\pmod7.
\end{equation*}
The field $L_0$ is the unramified cubic extension of $\mathbb Q_7$,
and logarithm identifies $U_0$ with $7\mathbb Z_7[\alpha]$.
The two faithful lines have bases $7\alpha$ and $7\alpha^2$.
Both coefficients in \eqref{seven:base-local-generator} are units.
As $\mathcal E_0^\chi$ has rank one, this proves
\begin{equation*}\tag{U5}\label{seven:base-local-isomorphism}
 \ell_0:\mathcal E_0^\chi\xrightarrow{\ \sim\ }U_0^\chi.
\end{equation*}

\subsection{Inverse limits and their control maps}

We use two elementary facts about compact modules. First, inverse
limit is exact on inverse systems of compact Hausdorff abelian groups.
To lift a compatible element of a quotient system, any finite set of
compatibility conditions can be met by lifting at its highest level
and taking images at lower levels. Compactness then supplies a lift
meeting all conditions. Transition maps need not be surjective.
In particular, for a compatible endomorphism $T$,
\begin{equation*}\tag{U6}\label{seven:compact-coinvariants}
 (\varprojlim B_n)/T(\varprojlim B_n)
 \simeq\varprojlim(B_n/TB_n).
\end{equation*}
Indeed, exactness applied first to $B_n\to TB_n$ proves
$\varprojlim TB_n=T\varprojlim B_n$, and a second application proves
\eqref{seven:compact-coinvariants}.

Second, the following compact form of Nakayama does not presuppose
finite generation. If a compact pro-seven $\Lambda$-module $B$ has
continuous finite $\Gamma$-stable quotients separating its points and
$B/(7,T)B$ is generated by the image of $b$, then $B=\Lambda b$.
The submodule $\Lambda b$ is compact and closed. Its quotient $Q$
satisfies $Q=(7,T)Q$. Every continuous finite quotient of $Q$ is a
module for some $\mathbb Z_7[C_{7^r}]$ on which $(7,T)$ is nilpotent,
and so must vanish. The finite-level unit groups modulo powers of
seven give the separating quotients needed here.

\begin{lemma}\label{seven:local-unit-control}
For $n>m$, norm induces an isomorphism
\[
 U_n^\chi/(\gamma^{7^m}-1)U_n^\chi
 \xrightarrow{\ \sim\ }U_m^\chi.
\]
Moreover $U_\infty^\chi\simeq\Lambda$, and
$U_\infty^\chi/\omega_mU_\infty^\chi\simeq U_m^\chi$.
\end{lemma}

\begin{proof}
The local extension $L_n/L_m$ is cyclic and totally ramified.
Local class field theory identifies
$U_m/N_{n/m}U_n$ with its Galois group. The $\Delta$-action on that
quotient is trivial because the unramified cubic and cyclotomic
ramified Galois groups commute. Hilbert 90 and the valuation sequence
also give
\[
 H^1(\operatorname{Gal}(L_n/L_m),\mathcal O_{L_n}^\times)
 \simeq\mathbb Z/7^{n-m}\mathbb Z
\]
with trivial $\Delta$-action. Replacing full units by principal units
does not change these Tate groups, since the residue unit group has
order $342$, prime to seven. Both
$\widehat H^0$ and $H^1$ therefore vanish on the faithful component.
The claimed finite norm isomorphism follows from the cyclic Tate
complex.

Compact exactness gives
$U_\infty^\chi/TU_\infty^\chi\simeq U_0^\chi\simeq\mathbb Z_7$.
Compact Nakayama makes $U_\infty^\chi$ cyclic over $\Lambda$.
All the finite local norms on this character are surjective, so a
generator of the limit maps to a generator of every $U_n^\chi$.
These groups are free of rank $7^n$ over $\mathbb Z_7$. There are no
seven-power roots of unity, since the local ramification index is a
power of seven, and their rational rank follows from the local
normal-basis representation and logarithm. Consequently the
surjection
\[
 \Lambda/\omega_n\Lambda\longrightarrow U_n^\chi
\]
is an isomorphism between free $\mathbb Z_7$-modules of the same rank.
The annihilator of a generator of $U_\infty^\chi$ lies in every
$\omega_n\Lambda$. But $\omega_n\in(7,T)^{n+1}$, and their intersection
is zero. This proves freeness. The same finite norm argument and
compact exactness give control at every level $m$.
\end{proof}

The global control map requires the capitulation terms, including
their transition maps.
For $G_n=\operatorname{Gal}(E_n/E)$ and $n\ge1$, the cyclic Tate
complex and Lemma~\ref{seven:unit-norm-index} give
\begin{equation*}\tag{U7}\label{seven:global-coinvariant-sequence}
 0\longrightarrow H^1(G_n,\mathcal E_n)^\chi
 \longrightarrow\mathcal E_n^\chi/T\mathcal E_n^\chi
 \xrightarrow{\ N_{n/0}\ }7\mathcal E_0^\chi
 \longrightarrow0.
\end{equation*}
The left term has order seven, but its norm transition from any higher
positive level to a lower one is zero. To see this, represent a base
class by an ideal $\mathfrak a$ and choose
$(z_n)=\mathfrak a\mathcal O_{E_n}$. Its cohomology class is represented
by $\gamma(z_n)/z_n$. For $n>m$, put $z_m'=N_{n/m}z_n$.
Then
\[
 (z_m')=\mathfrak a^{7^{n-m}}\mathcal O_{E_m}.
\]
The base class has order dividing seven, so
$\mathfrak a^{7^{n-m}}=(b)$ for some $b\in E^\times$.
Thus $z_m'/b$ is a unit, and
\[
 N_{n/m}\!\left(\frac{\gamma(z_n)}{z_n}\right)
 =\frac{\gamma(z_m'/b)}{z_m'/b}
\]
is a coboundary. This calculation commutes with the faithful
idempotent. The transition on the right of
\eqref{seven:global-coinvariant-sequence} is the identity.
Taking compact inverse limits therefore proves
\begin{equation*}\tag{U8}\label{seven:global-limit-coinvariants}
 \mathcal E_\infty^\chi/T\mathcal E_\infty^\chi
 \xrightarrow{\ \sim\ }7\mathcal E_0^\chi.
\end{equation*}
In particular there is no unaccounted kernel in this specialization.
Compact Nakayama now proves that $\mathcal E_\infty^\chi$ is cyclic.

The finite maps $\mathcal E_n\to\overline{\mathcal E}_n$ are surjective
maps of compact groups and commute with norms. Compact exactness
identifies $\overline{\mathcal E}_\infty$ with the image of
$\mathcal E_\infty\to U_\infty$, without requiring injectivity of
finite-level localization. Choose cyclic generators $e_\infty$ and
$u_\infty$ of $\mathcal E_\infty^\chi$ and $U_\infty^\chi$, and write
\[
 \ell_\infty(e_\infty)=f_{\mathfrak p,\chi}(T)u_\infty.
\]
Equations \eqref{seven:base-local-isomorphism} and
\eqref{seven:global-limit-coinvariants} give
\begin{equation*}\tag{U9}\label{seven:unit-factor-augmentation}
 f_{\mathfrak p,\chi}(0)\in7\mathbb Z_7^\times.
\end{equation*}
Thus $f_{\mathfrak p,\chi}\ne0$. Any relation on $e_\infty$ would
annihilate this nonzero element in the free local module $\Lambda$.
As $\Lambda$ is a domain, the global cyclic module is free and its
localization is injective. We have proved
\begin{equation*}\tag{U10}\label{seven:unit-factor}
 \mathcal E_\infty^\chi\simeq\Lambda,\qquad
 V_{\mathfrak p,\chi}\simeq\Lambda/(f_{\mathfrak p,\chi}).
\end{equation*}
The augmentation condition alone allows either $\mu=1,\lambda=0$
or $\mu=0,\lambda\ge1$. A first-layer calculation distinguishes them.

\subsection{An exact first-layer local index}

Use the period $\beta$, the automorphism $\tau$, and the thirteen
units of \S\ref{seven:finite-data}. Put $\pi=\beta+1$. The period
polynomial gives the Eisenstein equation
\begin{equation*}\tag{U11}\label{seven:local-eisenstein}
 \pi^7-7\pi^6+49\pi^4-98\pi^2-49\pi+7=0.
\end{equation*}
Thus $v_\pi(7)=7$, the residue field is
$\mathbb F_7[\overline\alpha]$ with $\overline\alpha^3=3$, and
$7/\pi^7\equiv-1\pmod\pi$. Set $U^j=1+\pi^j\mathcal O_{L_1}$.

\begin{lemma}\label{seven:local-finite-basis}
Let $\chi_a(\delta)=\zeta_3^a$, where
$\delta(\alpha)=\zeta_3\alpha$, $a=1,2$, and
$\zeta_3\equiv2\pmod7$. For
$J=\{1,2,3,4,5,6,8\}$, the seven elements
\[
 b_{a,j}=e_{\chi_a}(1+\alpha^a\pi^j),\qquad j\in J,
\]
form a $\mathbb Z_7$-basis of $U_1^{\chi_a}$.
Their images in $(U^1/U^{16})^{\chi_a}$ have respective orders
$343,49,49,49,49,49,49$ and form a direct-product basis.
\end{lemma}

\begin{proof}
For $i=0,1,2$, the successive leading grades of
$g_i=1+\alpha^i\pi$, $g_i^7$, and $g_i^{49}$ are $1,7,14$,
with coefficients
$\overline\alpha^i,\overline\alpha^{7i},-\overline\alpha^{7i}$.
For $h_{i,j}=1+\alpha^i\pi^j$, with
$j\in\{2,3,4,5,6,8\}$, the leading grades of $h_{i,j}$ and
$h_{i,j}^7$ are $j,j+7$, with coefficients
$\overline\alpha^i,-\overline\alpha^i$.
These assertions follow from the binomial expansion: the competing
lowest grades in $(1+c\pi^j)^7-1$ are $7+j$ and $7j$.
The next powers lie in $U^{16}$.

The displayed pivot grades partition $\{1,\ldots,15\}$, and the three
coefficients at each grade form an $\mathbb F_7$-basis of the residue
field. Successive digit elimination gives
\[
 U^1/U^{16}\simeq(C_{343})^3\oplus(C_{49})^{18}
\]
with these generators. On each associated graded quotient,
$e_{\chi_a}$ projects onto
$\mathbb F_7\overline\alpha^a$.
Moreover
$\overline\alpha^{7a}=3^{2a}\overline\alpha^a$ stays in that line.
The same elimination proves the asserted faithful direct-product
basis.

Logarithm and exponential are inverse on $U^j$ for $j\ge2$, since
$j>7/(7-1)$. Hence
$U^{16}=(U^9)^7\subseteq(U^1)^7$.
Reduction modulo $U^{16}$ therefore preserves
$U_1^{\chi_a}/7U_1^{\chi_a}$. The seven elements generate this quotient,
and $U_1^{\chi_a}$ is free of rank seven over $\mathbb Z_7$.
Nakayama proves that the displayed lifts form a full basis.
\end{proof}

The seven projected global units to compare with this basis are
\[
 e_{\chi_a}\varepsilon_0,\quad
 e_{\chi_a}\eta,\quad e_{\chi_a}\tau\eta,\quad\ldots,\quad
 e_{\chi_a}\tau^5\eta.
\]
They form a full $\mathbb Z_7$-basis of $\mathcal E_1^{\chi_a}$.
Indeed, prime-to-seven descent identifies
$\mathcal E(F_1)=\mathcal E(E_1)^{c=1}$, and projection gives an exact
sequence
\[
 0\longrightarrow\mathcal E(Q_1)
 \longrightarrow\mathcal E(F_1)
 \xrightarrow{\ e_{\chi_a}\ }\mathcal E(E_1)^{\chi_a}
 \longrightarrow0.
\]
For the kernel, the $\Delta$-trivial unit module is
$\mathcal E(KQ_1)$; its units are fixed by $c$ after seven-adic
completion. In fact $u/c(u)$ has absolute value one at all complex
embeddings and is a root of unity, and norm and extension identify
the completed units of this CM field with those of $Q_1$.
For surjectivity a vector $y$ in the faithful line lifts to $y+cy$.
The first six units in the thirteen-element seven-saturated basis
lie in the rank-six kernel. Their span is a direct summand, and hence
saturated, so it is the entire kernel. The remaining seven vectors
therefore give the stated basis.

Raise these seven global units to $342$ to make them principal.
This multiplies their completed coordinate columns by a
$\mathbb Z_7$-unit, and does not change their lattice index.
In the local basis of Lemma~\ref{seven:local-finite-basis}, the
coordinate matrices modulo $49$, with row order
$1,2,3,4,5,6,8$ and column order
$\varepsilon_0,\eta,\tau\eta,\ldots,\tau^5\eta$, are
\begin{equation*}\tag{U12}\label{seven:local-matrix-one}
 B_1=
 \begin{pmatrix}
 21&0&28&14&35&14&35\\
 0&48&48&34&48&27&48\\
 0&17&12&0&2&11&27\\
 28&45&24&25&41&2&34\\
 0&48&39&34&29&48&3\\
 7&26&15&30&16&22&12\\
 3&44&4&40&9&8&0
 \end{pmatrix},
 \qquad\det B_1\equiv28\pmod{49},
\end{equation*}
and
\begin{equation*}\tag{U13}\label{seven:local-matrix-two}
 B_2=
 \begin{pmatrix}
 42&42&21&0&14&21&28\\
 0&30&23&9&9&16&37\\
 0&20&44&33&29&32&28\\
 21&37&43&19&42&14&33\\
 0&22&32&31&13&14&28\\
 28&11&29&0&35&23&37\\
 13&44&0&12&34&23&32
 \end{pmatrix},
 \qquad\det B_2\equiv7\pmod{49}.
\end{equation*}

For precision, the finite arithmetic needed to verify these entries
can be performed entirely in successive $\pi$-adic digits. Equation
\eqref{seven:local-eisenstein} gives the exact carry rule
\begin{equation*}\tag{U14}\label{seven:local-carry}
 7\equiv-\pi^7-\pi^{13}+\pi^{15}\pmod{\pi^{16}}.
\end{equation*}
Use sixteen digits, each in the basis
$1,\alpha,\alpha^2$ with coefficients in $\{0,\ldots,6\}$.
Reduce $\alpha^3=269$ and carry every coefficient multiple of seven
upward by \eqref{seven:local-carry}; each carry strictly increases its
$\pi$-degree. This gives the ring
$\mathcal O_{L_1}/(\pi^{16})$ of order $7^{48}$.
All denominators in the printed unit expressions are prime to seven,
and their inverses modulo $343$ suffice, since $7^3=0$ in this ring.
Use $\beta=\pi-1$, the printed polynomial for $\tau$, and
$\zeta_3=324\pmod{343}$. On a principal unit the projector is the
product of its three $\Delta$-conjugates with exponents
$3^{-1}\zeta_3^{-ar}\pmod{343}$, $r=0,1,2$.
The pivot grades in the preceding lemma then recover the first
coordinate modulo $343$ and the other six modulo $49$.
This procedure determines all entries of
\eqref{seven:local-matrix-one}--\eqref{seven:local-matrix-two} from the
unit data already printed in \S\ref{seven:finite-data}.

Both determinants have valuation exactly one. The local and global
vectors used are full $\mathbb Z_7$-bases, so we have proved the exact
indices
\begin{equation*}\tag{U15}\label{seven:first-layer-local-index}
 [U_1^\chi:\overline{\mathcal E}_1^\chi]=7
\end{equation*}
for both faithful characters. The calculation applies at either
chosen prime above seven; changing the image of $\omega$ interchanges
the character labels.
As a further exact norm check on the same data, the printed expression
$\eta(\alpha,\beta)$ satisfies
\[
 \operatorname{Res}_B(q(B),\eta(\alpha,B))=1
 \quad\hbox{in }\mathbb Q[\alpha]/(\alpha^3-269).
\]
The six $Q_1$-units have relative norms
$-1,-1,1,-1,-1,1$, and $\varepsilon_0$ has relative norm
$\varepsilon_0^7$. This confirms the first-layer global norm defect
without serving as an additional premise for the all-level argument.

\subsection{The linear factor and the weight valuations}

For every $n>1$, Lemma~\ref{seven:unit-norm-index} gives
$[\mathcal E_1^\chi:N_{n/1}\mathcal E_n^\chi]=7$.
These norm images form a descending chain with the same index, hence
are equal. Compactness identifies the common image with the
projection of $\mathcal E_\infty^\chi$ to level one.
Lemma~\ref{seven:local-unit-control} therefore gives
\[
 V_{\mathfrak p,\chi}/\omega_1V_{\mathfrak p,\chi}
 \simeq
 U_1^\chi/
 \ell_1\bigl(\operatorname{im}
       (\mathcal E_\infty^\chi\to\mathcal E_1^\chi)\bigr).
\]
Equation \eqref{seven:first-layer-local-index} also proves that
$\ell_1$ is injective on this character, since its source and target
have equal rank. The universal-norm image has index seven in the
global units, so
\begin{equation*}\tag{U16}\label{seven:unit-factor-first-quotient}
 \#\Lambda/(f_{\mathfrak p,\chi},\omega_1)=7^2.
\end{equation*}

We now determine the Weierstrass invariants. If all coefficients of
$f=f_{\mathfrak p,\chi}$ are divisible by seven, then
\eqref{seven:unit-factor-augmentation} implies $f=7u$ for a unit $u$.
The quotient in \eqref{seven:unit-factor-first-quotient} would then
have order $7^7$, a contradiction. Otherwise Weierstrass preparation
writes $f=ug$, where $g$ is a distinguished polynomial of degree
$d\ge1$. Its constant coefficient has valuation one, so $g$ is
Eisenstein. If $1\le d\le5$, then at a primitive seventh root
$\zeta$ the leading term of $g(\zeta-1)$ has uniquely smallest
valuation $d/6$. Including the root $\zeta=1$, we obtain
\[
 v_7\!\left(\prod_{\zeta^7=1}g(\zeta-1)\right)=1+d.
\]
If $d\ge6$, each nontrivial-root factor has valuation at least one,
so the displayed valuation is at least seven, or the quotient is
infinite. Equation \eqref{seven:unit-factor-first-quotient} therefore
forces $d=1$. We have proved
\[
 f_{\mathfrak p,\chi}(T)
 =u(T)(T-a_{\mathfrak p,\chi}),\qquad
 a_{\mathfrak p,\chi}\in7\mathbb Z_7^\times.
\]

Rubin's class-field sequence
\cite[printed p.~36, equation~(5)]{Rubin} gives
\eqref{seven:one-prime-exact}. Its quotient is the finite faithful
class module of Theorem~\ref{seven:main}. Thus
$X_{\mathfrak p,\infty}^\chi$ is a finitely generated
$\mathbb Z_7$-module of rank one.
By \cite[Theorem~5.3(v), printed pp.~40--41]{Rubin},
$X_{\mathfrak p,\infty}$ has no nonzero pseudo-null submodule;
in the one-variable setting this excludes every finite submodule.
The $\mathbb Z_7$-torsion of $X_{\mathfrak p,\infty}^\chi$ is finite
and $\Lambda$-stable, so it vanishes. Hence this module is
$\mathbb Z_7$-free of rank one. Its $\Gamma$-action agrees, on the
finite-index submodule $V_{\mathfrak p,\chi}$, with multiplication
by $1+a_{\mathfrak p,\chi}$, and must agree on the whole module.
The submodule of index seven is consequently
$7X_{\mathfrak p,\infty}^\chi$. This proves all the asserted module
identifications.

The hypotheses of Rubin's theorem are satisfied: $K$ has class
number one and six roots of unity, $E/K$ is abelian of degree three,
and $E_\infty/K$ is abelian with
$\operatorname{Gal}(E_\infty/E)\simeq\mathbb Z_7$.
The statement concerns any such one-variable extension, so it
applies to the cyclotomic line; it is not restricted to the
$\mathfrak p$-ramified line.

It remains to evaluate the linear factor. Both
$\log(1+a_{\mathfrak p,\chi})$ and $\log8$ have valuation one, so
$e_{\mathfrak p,\chi}^{*}$ is a seven-adic unit and
$8^{-e_{\mathfrak p,\chi}^{*}}=1+a_{\mathfrak p,\chi}$.
For every $e\in\mathbb Z_7$,
\[
 8^{-e}-(1+a_{\mathfrak p,\chi})
 =(1+a_{\mathfrak p,\chi})
   \left(8^{-(e-e_{\mathfrak p,\chi}^{*})}-1\right).
\]
The logarithm and exponential on $7\mathbb Z_7$ give
$v_7(8^z-1)=1+v_7(z)$ for $z\in\mathbb Z_7$.
A unit in $\Lambda$ evaluates to a unit at every point of
$7\mathbb Z_7$. Equation
\eqref{seven:one-prime-weight-valuation} follows for every
characteristic generator.
Positive integers approximate any element of $\mathbb Z_7$ to
arbitrary precision, and hence these valuations are unbounded.

The standard Iwasawa involution replaces the root
$a_{\mathfrak p,\chi}$ by
$(1+a_{\mathfrak p,\chi})^{-1}-1$, which still has valuation one,
and replaces $e_{\mathfrak p,\chi}^{*}$ by its negative.
It does not remove the discrepancy with the constant strict length.
Thus identifying the strict length with the valuation of a generator
of this one-prime characteristic ideal is false, independently of
unit ambiguity and this involution.
An identification with a particular normalized Katz function still
requires its precise character, period, reciprocity, and Euler-factor
conventions. The conclusion here concerns the characteristic
generator of the explicitly defined one-prime module.
